\documentclass[12pt]{article}
\usepackage{fullpage}
\usepackage{amsmath,amssymb,amsthm}

\newtheorem{lemma}{Lemma}
\newtheorem{proposition}{Proposition}
\newtheorem{assumption}{Assumption}
\newcommand{\R}{\mathbb R}
\newcommand{\E}{\mathbb E}
\newcommand{\1}{\mathbf 1}
\newcommand{\cE}{\mathcal E}
\newcommand{\Var}{\operatorname{Var}}
\newcommand{\capacity}{\operatorname{Cap}}

\begin{document}

\begin{center}
{\bf Invariant measures for the Dirichlet skew stochastic heat equation}
\end{center}

\section*{Problem statement and interpretation}

Let \(E=C_0([0,1])\), and let \(S_t\) be the Dirichlet heat semigroup.
I use the ABLM regularized/mild interpretation specified in the problem,
including the assumed Dirichlet extension.  When passing from the
regularized equations to semigroups I use the following pointwise
interpretation of the stated weak convergence: for every bounded continuous
\(\Phi:E\to\R\), deterministic \(x\in E\), and \(t>0\),
\[
        P_t^n\Phi(x)\longrightarrow P_t\Phi(x).      \tag{A}
\]
The distributional drift has the gradient form
\[
        \delta_0=-\frac12\frac{d}{dr}(-2H_+),\qquad
        H_+(r)=\1_{\{r>0\}} .
\]
Thus the natural reversible measure is a bounded Gibbs perturbation of
Brownian bridge measure.  Let \(\gamma\) be the law on \(E\) of the
standard Brownian bridge from \(0\) to \(0\), and set
\[
   \Lambda(v)=\int_0^1\1_{\{v(x)>0\}}\,dx,\qquad
   \pi(dv)=Z^{-1}e^{2\Lambda(v)}\gamma(dv).       \tag{1}
\]
The map \(\Lambda\) is Borel.  Since \(0\le\Lambda\le1\), \(\pi\sim\gamma\)
with density bounded above and below.  The value assigned to \(H_+(0)\)
is irrelevant because a Brownian bridge spends zero Lebesgue time at
level \(0\):
\[
   \E_\gamma\int_0^1\1_{\{\beta(x)=0\}}\,dx
    =\int_0^1\mathbb P(\beta(x)=0)\,dx=0 .
\]

\section*{Consequences of the smooth approximations}

Put
\[
   B_n(r)=\int_{-\infty}^rG_{1/n}(s)\,ds,\qquad
   \pi_n(dv)=Z_n^{-1}\exp\!\left(2\int_0^1B_n(v(x))\,dx\right)\gamma(dv).
\]
Then \(0\le B_n\le1\), \(B_n(r)\to H_+(r)\) for \(r\ne0\), and hence
\[
       \|\pi_n-\pi\|_{\rm TV}\longrightarrow0 .          \tag{2}
\]
Indeed, the normalized densities converge in \(L^1(\gamma)\) by dominated
convergence and the zero-occupation identity above.

Let \(P_t^n\) be the semigroup of the regularized equation.  For each
fixed \(n\), the Nemytskii drift \(v\mapsto G_{1/n}(v)\) is bounded and
globally Lipschitz as an \(L^2(0,1)\)-valued map.  The standard
gradient-SPDE/Dirichlet-form theorem identifies its mild semigroup with
the symmetric diffusion generated, on smooth cylinder functions, by
\[
 L_nF(v)=\frac12{\rm Tr}\,D^2F(v)
       +\left\langle \frac12\Delta_Dv+G_{1/n}(v),DF(v)\right\rangle ,
\]
and by the closed form
\[
   \cE_n(F,G)=\frac12\int_E\langle DF,DG\rangle_{L^2(0,1)}\,d\pi_n .
\]
Equivalently,
\[
   \int G\,L_nF\,d\pi_n
   =-\frac12\int\langle DF,DG\rangle_{L^2(0,1)}\,d\pi_n
   =\int F\,L_nG\,d\pi_n.                         \tag{3}
\]
Thus \(P_t^n\) is reversible with invariant law \(\pi_n\).

\begin{proposition}\label{prop:limit}
Under (A), \(\pi\) is invariant for \(P_t\), \(P_t\) is symmetric on
\(L^2(\pi)\), and the \(L^2(\pi)\)-semigroup has a spectral gap.
\end{proposition}

\begin{proof}
For \(f,g\in C_b(E)\), dominated convergence with respect to \(\pi\),
(2), and invariance of \(\pi_n\) give
\[
 \int P_tf\,d\pi
 =\lim_n\int P_t^nf\,d\pi
  =\lim_n\int P_t^nf\,d\pi_n
  =\lim_n\int f\,d\pi_n
  =\int f\,d\pi .
\]
Similarly, using (3),
\[
 \int g\,P_tf\,d\pi
 =\lim_n\int g\,P_t^nf\,d\pi_n
  =\lim_n\int f\,P_t^ng\,d\pi_n
  =\int f\,P_tg\,d\pi .
\]
Since \(C_b(E)\) is dense in \(L^2(\pi)\), this extends to all of
\(L^2(\pi)\).

The Brownian bridge satisfies the Gaussian Poincare inequality
\[
     \Var_\gamma(F)\le C_0\int\|DF\|_{L^2(0,1)}^2\,d\gamma .
\]
Because the densities \(d\pi_n/d\gamma\) are bounded above and below
uniformly in \(n\), there is \(C<\infty\) such that
\[
     \Var_{\pi_n}(F)\le C\,\cE_n(F,F).             \tag{4}
\]
The spectral theorem for the reversible semigroup generated by
\(\cE_n\) therefore gives, with \(c>0\) independent of \(n\),
\[
     \Var_{\pi_n}(P_t^n f)
       \le e^{-2ct}\Var_{\pi_n}(f).                \tag{5}
\]
Passing to the limit for bounded continuous \(f\) and then using
\(L^2(\pi)\)-density yields
\[
     \Var_{\pi}(P_t f)\le e^{-2ct}\Var_\pi(f),
     \qquad f\in L^2(\pi).                         \tag{6}
\]
\end{proof}

\section*{Uniqueness once an all-starting-points upgrade is known}

\begin{lemma}\label{lem:criterion}
Assume Proposition \ref{prop:limit}.  If, for some \(\lambda>0\),
\[
     R_\lambda^P(x,\cdot)\ll \pi,\qquad
     R_\lambda^P(x,A):=\int_0^\infty e^{-\lambda t}P_t(x,A)\,dt,
     \quad x\in E,                                  \tag{7}
\]
then \(\pi\) is the only invariant probability measure of \(P_t\).
\end{lemma}

\begin{proof}
Let \(\eta\) be invariant.  If \(\pi(A)=0\), then (7) and invariance give
\[
   \eta(A)=\lambda\int_E R_\lambda^P(x,A)\,\eta(dx)=0.
\]
Thus \(\eta=h\pi\).  For bounded \(\varphi\),
\[
 \int \varphi h\,d\pi=\int P_t\varphi\,h\,d\pi
   =\int \varphi\,P_t h\,d\pi ,
\]
first for bounded \(h\), then by truncation.  Hence \(P_t h=h\) in
\(L^1(\pi)\).  With \(h_a=h\wedge a\), concavity and the Markov property
give \(h_a=(P_th)\wedge a\ge P_t h_a\).  The two sides have the same
\(\pi\)-integral, so equality holds.  Since \(h_a\in L^2(\pi)\), (6)
forces \(\Var_\pi(h_a)=0\).  Letting \(a\) vary shows \(h\equiv1\).
\end{proof}

\begin{lemma}\label{lem:sf}
Either of the following implies (7): \(P_t\) is strong Feller for every
\(t>0\), or, for one \(\lambda>0\), \(R_\lambda^P\) maps bounded Borel
functions to continuous functions.
\end{lemma}

\begin{proof}
If \(\pi(A)=0\), then \(x\mapsto P_t(x,A)\), respectively
\(x\mapsto R_\lambda^P(x,A)\), is a nonnegative continuous function with
zero \(\pi\)-integral.  Since \(\pi\sim\gamma\) and Brownian bridge
measure has full support on \(E\), that function is identically zero.
\end{proof}

A weaker topological replacement for strong Feller would also suffice.

\begin{lemma}[eventual equicontinuity criterion]\label{lem:eprop}
Assume Proposition \ref{prop:limit}.  Suppose that for every bounded
Lipschitz \(f:E\to\R\), every \(x\in E\), and every \(\varepsilon>0\),
there are a neighbourhood \(U\) of \(x\) and \(t_0>0\) such that
\[
        |P_t f(y)-P_t f(x)|\le\varepsilon,
        \qquad y\in U,\ t\ge t_0 .                  \tag{8}
\]
Then \(\pi\) is the unique invariant probability measure.
\end{lemma}

\begin{proof}
By (6), \(P_t f\to \pi(f)\) in \(L^2(\pi)\).  Fix \(x\) and choose \(U\)
as in (8).  Full support gives \(\pi(U)>0\).  For \(t\ge t_0\),
\[
 |P_t f(x)-\pi(f)|
 \le \varepsilon
 +\left(\pi(U)^{-1}\int_U |P_t f(y)-\pi(f)|^2\,\pi(dy)\right)^{1/2}.
\]
Letting \(t\to\infty\), then \(\varepsilon\downarrow0\), gives
\(P_t f(x)\to\pi(f)\) for every \(x\).  If \(\eta\) is invariant, dominated
convergence yields
\[
     \eta(f)=\int P_t f\,d\eta \longrightarrow \pi(f).
\]
Bounded Lipschitz functions determine probability measures on the Polish
space \(E\), so \(\eta=\pi\).
\end{proof}

\section*{The Dirichlet-form exceptional set}

For the bounded-variation potential \(f=-2H_+\), the Bounebache--Zambotti
Dirichlet form is
\[
        \cE(F,G)=\frac12\int_E\langle DF,DG\rangle_{L^2(0,1)}\,d\pi .
\]
Their Proposition 2.5 gives a resolvent kernel absolutely continuous with
respect to \(\pi\), but the representing identity is only \(\pi\)-a.e. in
the starting point; in the proof the kernel is extended over a
zero-capacity exceptional set.  Therefore it gives the needed resolvent
absolute continuity only away from a properly exceptional set \(N\).
Verification that the BZ local-time process lies in the Dirichlet ABLM
weak/mild class would identify the constructions at most quasi-everywhere;
this does not by itself settle deterministic \(x\in N\).

The following two reductions avoid asking for pointwise absolute
continuity on all of \(E\).

\begin{lemma}\label{lem:entrance}
Let \(N\) be a BZ properly exceptional set such that
\(R_\lambda^P(y,\cdot)\ll\pi\) for \(y\notin N\).  If
\[
          P_s(x,N)=0,\qquad x\in E,\ s>0,              \tag{9}
\]
then (7) holds for every \(x\in E\).
\end{lemma}

\begin{proof}
Let \(\pi(A)=0\).  For \(y\notin N\), \(R_\lambda^P(y,A)=0\).  For
\(\varepsilon>0\),
\[
 R_\lambda^P(x,A)
 \le \varepsilon
  +e^{-\lambda\varepsilon}\int_E R_\lambda^P(y,A)P_\varepsilon(x,dy)
 \le \varepsilon+\lambda^{-1}P_\varepsilon(x,N).
\]
Use (9) and let \(\varepsilon\downarrow0\).
\end{proof}

\begin{lemma}\label{lem:capacity}
Assume that every invariant probability \(\eta\) satisfies the finite
energy bound
\[
       \left|\int \varphi\,d\eta\right|
       \le C_\eta\,\cE_1(\varphi,\varphi)^{1/2}       \tag{10}
\]
for every bounded nonnegative element of \(D(\cE)\), represented by its
quasi-continuous version.
Then every invariant \(\eta\) charges no BZ exceptional set.  Consequently
the quasi-everywhere BZ resolvent absolute continuity implies
\(\eta\ll\pi\), and Lemma \ref{lem:criterion} gives uniqueness.
\end{lemma}

\begin{proof}
Let \(O\) be open and let \(e_O\) be an equilibrium potential with
\(e_O\ge1\) quasi-everywhere on \(O\) and
\(\cE_1(e_O,e_O)\le 2\capacity(O)\).  Applying (10) to bounded truncations of
the quasi-continuous version of \(e_O\) gives
\[
       \eta(O)\le \int e_O\,d\eta
          \le \sqrt2\, C_\eta\,\capacity(O)^{1/2}.
\]
Outer regularity of capacity implies that \(\eta(N)=0\) whenever
\(\capacity(N)=0\).  If \(\pi(A)=0\), invariance and the BZ resolvent identity
on \(N^c\) give
\[
   \eta(A)=\lambda\int_{N^c}R_\lambda^P(y,A)\,\eta(dy)=0.
\]
Thus \(\eta\ll\pi\), and the last step is Lemma \ref{lem:criterion}.
\end{proof}

\section*{Remaining open issues}

The argument proves uniqueness once any one of the all-starting-points
upgrades above is available.  What is still unproved from the hypotheses
stated in the problem is an upgrade excluding invariant probabilities
living on the \(\pi\)-null exceptional set.  Concretely, one still needs
one of the following.

\begin{enumerate}
\item Prove (7), or prove one of the strong-Feller sufficient conditions in
Lemma \ref{lem:sf}.
\item Prove the eventual equicontinuity estimate (8).  The \(L^2(\pi)\)
spectral gap plus full support would then imply pointwise convergence
\(P_t f(x)\to\pi(f)\) by Lemma \ref{lem:eprop}.
\item Prove entrance into the BZ non-exceptional set, namely (9).
\item Prove the finite-energy/capacity estimate (10) for every stationary
law.  This would show that invariant laws are smooth measures for the BZ
Dirichlet form and hence cannot charge exceptional sets.
\end{enumerate}

This round I checked two non-pointwise routes: the eventual
equicontinuity route, which would upgrade \(L^2(\pi)\) mixing to pointwise
mixing, and the finite-energy route, which would only need invariant laws
to ignore capacity-zero sets.  Both reductions are rigorous as stated, but
the required estimates are not consequences of the assumptions currently
listed in the problem.  Therefore the present document proves uniqueness
among invariant measures absolutely continuous with respect to \(\pi\) and
gives several exact sufficient criteria for full uniqueness, but it does
not yet constitute an unconditional proof or disproof of the statement.

\begin{thebibliography}{9}
\bibitem{ABLM}
S. Athreya, O. Butkovsky, K. L\^e and L. Mytnik,
\emph{Well-posedness of stochastic heat equation with distributional drift
and skew stochastic heat equation}, Comm. Pure Appl. Math. 77 (2024),
2708--2777.

\bibitem{ABLMweakmild}
S. Athreya, O. Butkovsky, K. L\^e and L. Mytnik,
\emph{Analytically weak and mild solutions to stochastic heat equation with
irregular drift}, Stoch. PDE: Anal. Comp. (2025).

\bibitem{BZ}
S. K. Bounebache and L. Zambotti,
\emph{A skew stochastic heat equation}, J. Theoret. Probab. 27 (2014),
168--201.
\end{thebibliography}

\end{document}
